\subsection{Métodos}

\subsubsection{Regla: PascalCase y verbo descriptivo}

\textbf{Regla:} El nombre de un método deberá escribirse en \texttt{PascalCase} y comenzar con un verbo o frase verbal que describa la acción que realiza.

\textbf{Descripción:} Microsoft establece \texttt{PascalCase} como la convención para todos los métodos públicos y privados en C\#, a diferencia de otros lenguajes que usan minúscula inicial (Microsoft, 2025c).

\textbf{Código correcto:}
\begin{lstlisting}[style=csharpstyle]
public int CalculateTotalScore(int basePoints, int bonusPoints)
{
    return basePoints + bonusPoints;
}
\end{lstlisting}

\textbf{Código incorrecto:}
\begin{lstlisting}[style=csharpstyle]
public int totalScore(int basePoints, int bonusPoints)
{
    return basePoints + bonusPoints;
}
\end{lstlisting}

\subsubsection{Regla: máximo de parámetros}

\textbf{Regla:} Un método no deberá recibir más de tres parámetros. Cuando la operación requiera más datos de entrada, estos deberán encapsularse en un objeto de transferencia de datos (DTO).

\textbf{Descripción:} McConnell recomienda limitar el número de parámetros de una rutina, ya que una lista larga de parámetros incrementa la probabilidad de invocarla con el orden o el valor incorrecto (McConnell, 2004).

\textbf{Código correcto:}
\begin{lstlisting}[style=csharpstyle]
public readonly record struct MoveRequest(Player Player, Tile Origin, Tile Destination);

public sealed class MoveExecutor
{
    public MoveResult ExecuteMove(MoveRequest moveRequest)
    {
        return new MoveResult(moveRequest.Player, moveRequest.Destination);
    }
}
\end{lstlisting}

\textbf{Código incorrecto:}
\begin{lstlisting}[style=csharpstyle]
public sealed class MoveExecutor
{
    public MoveResult ExecuteMove(Player player, Tile origin, Tile destination, int diceRoll, bool isDoubleTurn)
    {
        return new MoveResult(player, destination);
    }
}
\end{lstlisting}

\subsubsection{Regla: métodos booleanos como pregunta}

\textbf{Regla:} Un método que devuelva \texttt{bool} deberá nombrarse con un prefijo interrogativo (\texttt{Is}, \texttt{Has}, \texttt{Can}) seguido del predicado que evalúa.

\textbf{Descripción:} Esta convención extiende a los métodos la misma recomendación que Martin aplica a las variables booleanas: el nombre debe leerse como la pregunta que el método responde (Martin, 2008).

\textbf{Código correcto:}
\begin{lstlisting}[style=csharpstyle]
public bool IsMoveValid(Tile origin, Tile destination)
{
    return destination.IsWithinBoard && !destination.IsOccupied;
}
\end{lstlisting}

\textbf{Código incorrecto:}
\begin{lstlisting}[style=csharpstyle]
public bool CheckMove(Tile origin, Tile destination)
{
    return destination.IsWithinBoard && !destination.IsOccupied;
}
\end{lstlisting}
